\color{unidarkgreen}{\subsection{Защиты по напряжению}}
\color{black}

\begin{enumerate}

\item Общие сведения
\begin{enumerate}

\item 
Функциональный блок ЗН БСК включает в себя две ступени защиты от повышения напряжения и одну ступень защиты минимального напряжения.

\item 
Уставки ЗН БСК приведены в таблице \ref{app_set:zn}.

\end{enumerate}

\item Защита от повышения напряжения 
\begin{enumerate}

\item
Защита от повышения напряжения предназначена как для защиты БСК от перегрузки по напряжению, так и для предотвращения недопустимых рабочих режимов сети с возможным нарушением устойчивости. 

\item
Защита содержит сигнальную и отключающие ступени. Каждая из ступеней содержит три реле напряжения максимального действия. В зависимости от положения программной накладки \textbf{«ЗПН измер.»} ступени могут контролировать междуфазные (Uab, Ubc, Uca) или фазные (Ua, Ub, Uc) напряжения. 

\item
Пуск ступени может быть произведен при повышении всех трех (схема «И») или повышении любого из контролируемых напряжений (схема «ИЛИ»). Выбор режима осуществляется программной накладкой \textbf{«Режим ЗПН»}. 

\item
Каждая из ступеней имеет регулируемое значение по напряжению срабатывания, задаваемое параметром \textbf{«Uсигн ЗПН»} и \textbf{«Uоткл ЗПН»} для сигнальной и отключающей ступени соответственно. 

\item
Сигнальная ступень с задержкой по времени \textbf{«Тсигн ЗПН»} действует на сигнализацию, отключающая ступень с задержкой по времени \textbf{«Тоткл ЗПН»} действует на отключение выключателя БСК. Алгоритм работы сигнальной ступени ЗПН приведен на рисунке \ref{pict:ZPN}, алгоритм работы отключающей ступени аналогичен алгоритму работы сигнальной ступени.

\medskip
\begin{figure}[H]
\centering
\includegraphics[width=1\textwidth,height=1\textheight,keepaspectratio]{\fbpath/09.01 БСК/261.0861 ЗН/img/ZPN.pdf}
\captionof{figure}{Функциональнальная схема логики ЗПН}\label{pict:ZPN}
\end{figure}

\end{enumerate}

\item Защита минимального напряжения 
\begin{enumerate}

\item
Защита минимального напряжения реагирует на снижение напряжения на шинах и предназначена для предотвращения повреждения БСК при исчезновении напряжения на шинах и повторной подачи напряжения на шины без предварительного отключения БСК (в этом случае к БСК может быть приложено удвоенное напряжение).

\item
Защита содержит три реле напряжения минимального действия, включенных по схеме «И». В зависимости от положения программной накладки \textbf{«ЗМН измер.»} ступени могут контролировать междуфазные (Uab, Ubc, Uca) или фазные (Ua, Ub, Uc) напряжения.

\item
Для исключения излишнего пуска защиты предусмотрено ее блокирование при отключенном состоянии выключателя. Ввод контроля положения выключателя в логику работы защиты осуществляется программной накладкой \textbf{«ЗМН конт.выкл.»}. 

\item
Защита имеет регулируемые значения по напряжению и времени срабатывания, задаваемые параметрами \textbf{«Uср ЗМН»} и \textbf{«Тср ЗМН»} соответственно.  Алгоритм работы ЗМН приведен на рисунке \ref{pict:ZMN}

\medskip
\begin{figure}[H]
\centering
\includegraphics[width=1\textwidth,height=1\textheight,keepaspectratio]{\fbpath/09.01 БСК/261.0861 ЗН/img/ZMN.pdf}
\captionof{figure}{Функциональнальная схема логики ЗМН}\label{pict:ZMN}
\end{figure}

\end{enumerate}


\end{enumerate} 